% Macros - Jia-Bin Huang (jbhuang@vt.edu)

\renewcommand{\baselinestretch}{1}       % for squeezing the draft into the page limit, do not use

% =========================================
% Useful macros
% =========================================

% Latin abbreviations
\def\etal{et al.~}			  % and others, and co-workers
% \newcommand\etal{et al.~}
\def\eg{e.g.,~}               % for example
\def\ie{i.e.,~}               % that is, in other words
\def\etc{etc}                 % and other things, and so forth
\def\cf{cf.~}                 % compare
\def\viz{viz.~}               % namely, precisely
\def\vs{vs.~}                 % against


% Math related
\DeclareMathOperator*{\argmin}{\arg\!\min} 
\DeclareMathOperator*{\argmax}{\arg\!\max}

% Consistent margin adjustment for paragraphs, figures, and sections
\newlength\paramargin
\newlength\figmargin
\newlength\secmargin
\newlength\figcapmargin

\setlength{\secmargin}{0.0mm}
\setlength{\paramargin}{0.0mm}
\setlength{\figmargin}{0.0mm}
\setlength{\figcapmargin}{0.5mm}

\setlength{\parskip}{0pt}     % 段落間距 0
\setlength{\parindent}{1em}   % 段落開頭縮排（1em 可調）

\newcommand{\red}{\textcolor{red}}
\newcommand{\blue}{\textcolor{blue}}

% minipage
\newcommand{\mpage}[2]
{
\begin{minipage}{#1\linewidth}\centering
#2
\end{minipage}
}

\newcommand{\mfigure}[2]
{
\begin{subfigure}[b]{#1\linewidth}\centering
\includegraphics[width=\linewidth]{#2}
\end{subfigure}
}


\newcommand{\Paragraph}[1]
{
\vspace{\paramargin}
\paragraph{#1}
}

\newcommand{\heading}[1]
{
\vspace{1mm}
\noindent \textbf{#1}
}   

% References for figures, tables, equations, and sections
\newcommand{\secref}[1]{Section~\ref{sec:#1}}
\newcommand{\figref}[1]{Figure~\ref{fig:#1}} 
\newcommand{\tabref}[1]{Table~\ref{tab:#1}}
\newcommand{\eqnref}[1]{Equation~\ref{eq:#1}}
\newcommand{\thmref}[1]{Theorem~\ref{#1}}
\newcommand{\prgref}[1]{Program~\ref{#1}}
\newcommand{\algref}[1]{Algorithm~\ref{#1}}
\newcommand{\clmref}[1]{Claim~\ref{#1}}
\newcommand{\lemref}[1]{Lemma~\ref{#1}}
\newcommand{\ptyref}[1]{Property~\ref{#1}}

% Comments
\long\def\ignorethis#1{}
\newcommand {\jiabin}[1]{{\color{blue}\textbf{Jia-Bin: }#1}\normalfont}
\newcommand {\james}[1]{{\color{violet}\textbf{James: }#1}\normalfont}
\newcommand {\eric}[1]{{\color{magenta}\textbf{Eric: }#1}\normalfont}
\newcommand {\todo}{{\textbf{\color{red}[TO-DO]\_}}}
\def\newtext#1{\textcolor{blue}{#1}}
\def\modtext#1{\textcolor{red}{#1}}
\newcommand{\note}[1]{{\it\color{blue} #1}}

\newcommand{\tb}[1]{\textbf{#1}}
\newcommand{\mb}[1]{\mathbf{#1}}

\newcommand{\jbox}[2]{
  \fbox{%
  	\begin{minipage}{#1}%
  		\hfill\vspace{#2}%
  	\end{minipage}%
  }}

\newcommand{\jblock}[2]{%
	\begin{minipage}[t]{#1}\vspace{0cm}\centering%
	#2%
	\end{minipage}%
}


% \titleformat{\section}
%   {\normalfont\bfseries\centering}
%   {\thesection.}{1em}{}

% % 將 \subsection 調成 10pt 加粗
% \titleformat{\subsection}
%   {\normalsize\bfseries} % 10pt = \normalsize
%   {\thesubsection.}{1em}{}

% % 格式設定：斜體、不加粗、10pt、編號、左對齊
% \titleformat{\subsubsection}
%   {\normalsize\itshape}      % 10pt, 斜體
%   {\thesubsubsection.}       % 顯示 5.1.1. 編號
%   {1em}                      % 編號與文字之間距
%   {}                         % 不加其他

% % 間距設定：上間距1ex，下間距0.3ex（很緊湊）
% \titlespacing*{\subsubsection}
%   {0pt}       % 左邊距
%   {1.5ex plus .3ex minus .3ex}  % 上方距離
%   {0.3ex}     % 下方距離（正文與標題之間）

% --- Section 樣式設定 ---
\titleformat{\section}
  {\normalsize\bfseries\centering} % 10pt 粗體 置中
  {\thesection.}{0.25em}{}

\titlespacing*{\section}
  {0pt}{1.2em}{0.8em} % 上下空白（可調整為更精確 10pt 單位）

% --- Subsection 樣式設定 ---
\titleformat{\subsection}
  {\normalsize\bfseries}  % 10pt 粗體 左對齊
  {\thesubsection.}{0.25em}{}

\titlespacing*{\subsection}
  {0pt}{1.2em}{0.5em}

% --- Subsubsection 樣式設定 ---
\titleformat{\subsubsection}
  {\normalsize\itshape}  % 10pt 斜體 左對齊
  {\thesubsubsection.}{0.25em}{}

\titlespacing*{\subsubsection}
  {0pt}{1.2em}{0.3em}